\subsection{不可约表示的存在性}

沿用上一节的记号。

定理 3（赖科夫）。--- Pos1(G) 上的弱拓扑与紧收敛拓扑重合。

先证明几个引理。

引理 5. --- 设 X 是局部紧拓扑空间，\(\nu\) 是 X 上的正测度。在 \(\mathcal{C}_{b}(X)\) 的每个有界子集上，由弱拓扑 \(\sigma(L^{\infty}(X), L^{1}(X))\) 诱导的拓扑比紧收敛拓扑粗。

设 B 是 \(\mathcal{C}_{b}(\mathrm{X})\) 的有界子集，且 \(M \in R_{+}\) 满足 \(\|\varphi\|_{\infty} \leqslant M\) 对所有 \(\varphi \in B\)。设 \(\psi \in L^{1}(\mathrm{X}, \nu)\) 且 \(\varepsilon > 0\)。取 X 的一个紧子集 K，使得

\[
\int_ {\mathrm{X-K}} | \psi | d \nu <   \varepsilon .
\]

于是对 B 中任意 \(\varphi_{1}\) 和 \(\varphi_{2}\)，有

\[
\begin{array}{r l} & {\langle \varphi_ {1} - \varphi_ {2}, \psi \rangle = \int_ {\mathrm{X}} \psi (\varphi_ {1} - \varphi_ {2}) d \nu} \\ & {\qquad = \int_ {\mathrm{K}} \psi (\varphi_ {1} - \varphi_ {2}) d \nu + \int_ {\mathrm{X-K}} \psi (\varphi_ {1} - \varphi_ {2}) d \nu ,} \end{array}
\]

从而

\[
| \langle \varphi_ {1} - \varphi_ {2}, \psi \rangle | \leqslant \sup _ {x \in \mathrm{K}} | \varphi_ {1} (x) - \varphi_ {2} (x) | \| \psi \| _ {1} + 2 \mathrm{M} \varepsilon ,
\]

引理得证。

引理 6. --- 设 \(\psi \in \mathrm{L}^{1}(\mathrm{G})\)。令 B 是巴拿赫空间 \(\mathrm{L}^{\infty}(\mathrm{G})\) 的有界子集。当 B 赋予弱拓扑而映射 \(\varphi \mapsto \psi * \varphi\) 的值域 \(\mathcal{C}_b(\mathrm{G})\) 赋予紧收敛拓扑时，该映射从 B 到 \(\mathcal{C}_b(\mathrm{G})\) 连续。

设 \(\varphi\in\mathrm{L}^{\infty}(\mathrm{G})\) 且 \(\eta\in\mathrm{L}^{1}(\mathrm{G})\)。函数 \(\Delta^{-1}\check{\eta}\) 属于 \(\mathrm{L}^{1}(\mathrm{G})\)，且 \(\langle\eta,\check{\psi}\rangle=\langle\Delta^{-1}\check{\eta},\psi\rangle\)（参见 INT, VII, p. 19, § 1, n°3, 公式 (22)）。因此，映射 \(\varphi\mapsto\check{\varphi}\) 是赋予弱拓扑的空间 \(\mathrm{L}^{\infty}(\mathrm{G})\) 上的自同构。

设 \(\varphi\in\mathrm{L}^{\infty}(\mathrm{G})\)。由 INT, VIII, p. 167, § 4, no. 5, 命题 14，函数 \(\psi*\varphi\) 属于 \(\mathcal{C}_{b}(\mathrm{G})\)，且满足

\[
\begin{array}{r l} & {(\psi * \varphi) (g) = \int_ {\mathrm{G}} \psi (h) \varphi (h ^ {- 1} g) d \mu (h)} \\ & {\qquad = \int_ {\mathrm{G}} \psi (g y) \check {\varphi} (y) d \mu (y) = \langle \check {\varphi}, \pmb {\gamma} _ {\mathrm{G}} ^ {(1)} (g ^ {- 1}) \psi \rangle} \end{array}
\]

对所有 \(g \in G\) 都成立。因此，线性映射 \(u: \varphi \mapsto \psi * \varphi\) 从赋予弱拓扑的空间 \(L^{\infty}(G)\) 到赋予逐点收敛拓扑的 \(\mathcal{C}_{b}(G)\) 是连续映射。

设 \(\varphi\in\mathrm{B}\) 且 \((g,h)\in\mathrm{G}\times\mathrm{G}\)。由上式，有

\[
| u (\varphi) (g) - u (\varphi) (h) | \leqslant \| \check {\varphi} \| _ {\infty} \| (\boldsymbol {\gamma} _ {\mathrm{G}} ^ {(1)} (g ^ {- 1}) - \boldsymbol {\gamma} _ {\mathrm{G}} ^ {(1)} (h ^ {- 1})) \psi \| _ {1}.
\]

由于 B 有界且 G 在 \(L^{1}(G)\) 上的左正则表示连续（V, p. 405, no. 4），这意味着 \(u(B)\) 是 \(\mathcal{C}_{b}(G)\) 的等度连续子集。由前述结果及 TG, X, p. 16, 定理 1，断言得证。

引理 7. --- 设 \(\psi \in \mathrm{L}^1 (\mathbf{G})\) 满足 \(\psi \geqslant 0\) 且 \(\int \psi = 1\)。令 p 表示 \(\mathcal{C}_b(\mathbf{G})\) 上由 \(p(\varphi) = |\langle \varphi ,\psi \rangle |\) 对所有 \(\varphi \in \mathcal{C}_b(\mathbf{G})\) 定义的半范数。对每个 \(\varphi \in \mathrm{Pos}_1(\mathbf{G})\)，有

\[
\| \psi * \varphi - \varphi \| _ {\infty} \leqslant \sqrt {2 p (1 - \varphi)}.
\]

设 \(\varphi\in\mathrm{Pos}_{1}(\mathrm{G})\) 且 \(x\in\mathrm{G}\)。根据 INT, VIII, p. 167, § 4, no. 5, 命题 14，得到

\[
\begin{array}{r l} | \psi * \varphi (x) - \varphi (x) | = & \left| \int_ {\mathrm{G}} (\varphi (y ^ {- 1} x) - \varphi (x)) \psi (y) d \mu (y) \right| \\ & \leqslant \int_ {\mathrm{G}} | \varphi (y ^ {- 1} x) - \varphi (x) | \psi (y) d \mu (y) \\ & \leqslant \int_ {\mathrm{G}} \sqrt {2 (1 - \mathcal {R}   \varphi (y))} \psi (y) d \mu (y) \end{array}
\]

应用 V, p. 446 的界 (1)。于是由柯西–施瓦茨不等式可得

\[
\begin{array}{c} \| \psi * \varphi - \varphi \| _ {\infty} \leqslant \sqrt {2} \Bigl (\int_ {\mathrm{G}} (1 - \mathcal {R} (\varphi)) \psi   d \mu \Bigr) ^ {1 / 2} \Bigl (\int_ {\mathrm{G}} \psi   d \mu \Bigr) ^ {1 / 2} \\ \leqslant \sqrt {2} \sqrt {p (1 - \varphi)}, \end{array}
\]

从而得到结论。

引理 8. --- 设 K 是拓扑域，E 和 F 是 K 上的拓扑向量空间。设 f 是从 E 到 F 的映射，X 是 E 的子集，且 \(x \in X\)。若对于 F 中 0 的每个邻域 W，都存在 E 中 x 的一个邻域 U 以及从 X 到 F 的连续映射 g，使得映射 f \textbar{} X 在 x 处连续，当 \((f - g)(\mathrm{U} \cap \mathrm{X}) \subset \mathrm{W}\)。

设 \(W_{0}\) 是 F 中 0 的一个邻域，\(V_{0}\) 是 F 中 0 的一个对称邻域，满足 \(V_{0} + V_{0} + V_{0} \subset W_{0}\)。设 \(U_{0}\) 是 E 中 x 的一个邻域，g 是从 X 到 F 的连续映射，满足 \((f - g)(\mathrm{U}_{0} \cap \mathrm{X}) \subset \mathrm{V}_{0}\)。

存在 E 中 x 的一个邻域 \(U_{1}\)，使得 \(g(\mathrm{U}_{1} \cap \mathrm{X}) \subset g(x) + \mathrm{V}_{0}\)。对每个 \(y \in U_{0} \cap U_{1} \cap X\)，有

\[
\begin{array}{c} f (y) - f (x) = (f (y) - g (y)) + (g (y) - g (x)) + (g (x) - f (x)) \in V_{0} + V_{0} + V_{0} \subset W_{0}, \end{array}
\]

所以 \(f(y) \in f(x) + W_{0}\)，从而得到结论。

现在证明定理 3。以 \(\mathrm{Pos}_{1}(\mathrm{G})_{f}\)（分别以 \(\mathrm{Pos}_{1}(\mathrm{G})_{c}\)）表示赋予弱拓扑（分别赋予紧收敛拓扑）的集合 \(\mathrm{Pos}_{1}(\mathrm{G})\)；同样，以 \(\mathcal{C}_{b}(\mathrm{G})_{f}\)（分别以 \(\mathcal{C}_{b}(\mathrm{G})_{c}\)）表示赋予弱拓扑（分别赋予紧收敛拓扑）的空间 \(\mathcal{C}_{b}(\mathrm{G})\)。

从 \(\mathrm{Pos}_{1}(\mathrm{G})_{c}\) 到 \(\mathrm{Pos}_{1}(\mathrm{G})_{f}\) 的恒等映射连续（引理 5）。反之，以 \(\iota\) 表示 \(\mathrm{Pos}_{1}(\mathrm{G})_{f}\) 到 \(\mathcal{C}_{b}(\mathrm{G})_{c}\) 的包含映射。证明 \(\iota\) 连续即可完成证明。

设 \(\varphi_{1} \in \mathrm{Pos}_{1}(\mathbf{G})\)。为验证 \(\iota\) 在 \(\varphi_{1}\) 处连续，应用引理 8。设 W 是 \(\mathcal{C}_{b}(\mathbf{G})_{c}\) 中 0 的一个邻域。只需找到线性映射 \(u\)，从 \(\mathcal{C}_{b}(\mathbf{G})\) 到 \(\mathcal{C}_{b}(\mathbf{G})\)，它通过对子空间的传递诱导出从 \(\mathrm{Pos}_{1}(\mathbf{G})_{f}\) 到 \(\mathcal{C}_{b}(\mathbf{G})_{c}\) 的连续映射，并找到 \(\mathrm{Pos}_{1}(\mathbf{G})_{f}\) 中 0 的一个邻域 U，使得 \((\iota - u)(\mathrm{U} \cap \mathrm{Pos}_{1}(\mathbf{G})) \subset \mathrm{W}\)。

存在 \(\varepsilon > 0\)，使得 W 包含满足 \(\varphi \in \mathcal{C}_b(\mathrm{G})\) 且 \(\| \varphi \|_{\infty} \leqslant \varepsilon\) 的函数所成的集合。由于 \(\varphi_1(e) = 1\)，存在 G 中 \(e\) 的一个紧邻域 V，使得

\[
\sup _ {x \in \mathrm{V}} | 1 - \varphi_ {1} (x) | \leqslant \frac {\varepsilon^ {2}}{4}.
\]

令 \(\varphi_{\mathrm{V}}\) 为 V 的特征函数，并置 \(\psi_{\mathrm{V}} = \mu (\mathrm{V})^{-1}\varphi_{\mathrm{V}}\)。线性映射 \(u\colon \varphi \mapsto \psi_{\mathrm{V}}*\varphi\) 从 \(\mathrm{Pos}_1(\mathrm{G})_f\) 到 \(\mathcal{C}_b(\mathrm{G})_c\) 连续（引理 6）。

以 \(q_{V}\) 表示由 \(\varphi\mapsto|\langle\varphi,\psi_{V}\rangle|\) 在 \(\mathcal{C}_{b}(G)\) 上定义的半范数；它对于弱拓扑连续。由于 \(\psi_{V}\) 在 V 外为零，可得 \(q_{\mathrm{V}}(1-\varphi_{1})\leqslant\varepsilon^{2}/4\)；因此在 \(\varphi_{1}\) 的 \(\mathcal{C}_{b}(\mathrm{G})_{f}\) 中存在一个邻域 U，使得 \(q_{\mathrm{V}}(1-\varphi)\leqslant\varepsilon^{2}/2\) 对所有 \(\varphi\in U\) 成立。

设 \(\varphi\in\mathrm{U}\cap\mathrm{Pos}_{1}(\mathrm{G})\)。由引理 7，有

\[
\| (\iota - u) (\varphi) \| _ {\infty} \leqslant \sqrt {2 q _ {\mathrm{V}} (1 - \varphi)} \leqslant \varepsilon
\]

所以 \((\iota - u)(\mathrm{U}) \subset \mathrm{W}\)。定理得证。

一般而言，\(\operatorname{Pos}_{\leqslant1}(\mathbf{G})\) 上的弱拓扑不与紧收敛拓扑重合，群 R 的例子说明了这一点。

回忆我们以 \(\widehat{G}\) 表示 \(G\) 的不可约酉表示的等价类集合（V, p. 393, 定义 8）。

定理 4（盖尔范德--赖科夫）。--- 对 G 中每个 \(x \neq e\)，存在 G 的不可约酉表示 \(\pi\)，使得 \(\pi(x)\) 不是 \(\pi\) 的空间上的恒等自同态。

设 \(x \in G\) 满足 \(\pi(x)\) 对每个 \(\pi \in \widehat{G}\) 都是恒等映射。则有 \(\varphi(x) = \varphi(e) = 1\) 对每个极点 \(\varphi\)（属于 \(\mathrm{Pos}_{1}(\mathrm{G})\)）（V, p. 444, 命题 11）；又由于命题 14，这对每个 \(\varphi \in \mathrm{Pos}_{1}(\mathrm{G})\) 都成立，因为在线性泛函 \(\varphi \mapsto \varphi(e)\) 在赋予弱拓扑的 \(\mathrm{Pos}_{1}(\mathrm{G})\) 上是连续的（定理 3）。但是若 \(x \neq e\)，则存在范数为 1 的 \(f \in L^{2}(\mathrm{G})\)，使得 \(\langle f | \boldsymbol{\gamma}_{\mathrm{G}}(x)f \rangle = 0\)（V, p. 406, 引理 3）。由于这个等式左端具有 \(\varphi(x)\) 的形式，其中 \(\varphi \in \mathrm{Pos}_{1}(\mathrm{G})\)，故矛盾。

若 G 不是局部紧的，则 G 的不可约酉表示未必能分离 G 的点。

推论。--- 对 G 中任意 \(g \neq h\)，存在不可约表示 \(\pi \in \widehat{G}\) 及 \(\pi\) 的一个矩阵系数 f，使得 \(f(g) \neq f(h)\)。特别地，\(\Upsilon(\mathrm{G})\) 是 \(\mathcal{C}_{b}(\mathrm{G})\) 的能分离 G 的点的子代数。

设 \(g \neq h\) 属于 G。存在不可约表示 \(\pi \in \widehat{G}\)，使得 \(\pi(g) \neq \pi(h)\)（定理 4），所以在 \(\pi\) 的空间中存在 x 和 y，使得 \(\langle x | \pi(g)y \rangle \neq \langle x | \pi(h)y \rangle\)。

